\section{中层解析 I：cli（CLI、命令分发、业务编排）}

\subsection{模块定位}

\fpath{cli/} 是“用户入口层 + 中层业务编排层”：

\begin{itemize}
  \item \textbf{入口}：\fpath{cli/app/demo_main.cc} 解析 CLI 并决定执行模式；
  \item \textbf{命令分发}：\fpath{cli/shell/dispatch/router/dispatch.cc}、\fpath{cli/shell/dispatch/router/dispatch_routing.cc}（phase-2 前缀快路径）、\fpath{cli/shell/dispatch/registry/*}（命令表与快照源）与 \fpath{cli/shell/dispatch/handlers/*} 负责交互命令解析与路由；解析/索引胶水在 \fpath{cli/shell/dispatch/support/}，跨 handler 内部声明在 \fpath{cli/shell/dispatch/shared/}。
  \item \textbf{查询中间层}：\fpath{cli/modules/where/executor/*} 管理条件解析、候选集构建、策略限流、查询缓存与 \fpath{table_stats}（全表扫描构建统计；\texttt{*.tablestats} 文件持久化由环境变量控制）；\texttt{SHOW PLAN}/\texttt{EXPLAIN WHERE} 输出计划 JSON 或文本摘要。
  \item \textbf{事务协调}：\fpath{cli/modules/txn/coordinator/*} 封装 Txn + WAL + 文件锁 + vacuum 触发；\fpath{stats_impl.cc} 聚合 \texttt{TxnRuntimeStats}（含 \texttt{table\_storage\_health\_*}），供 \texttt{SHOW TUNING}/\texttt{SHOW TUNING JSON} 与 GUI 观测面板消费。
  \item \textbf{运行编排}：\fpath{cli/shell/bootstrap/demo_runner.cc} 对接 batch/exec/import/mdb 等模式。
\end{itemize}

\subsection{关键数据结构（中层）}

\subsubsection{ShellState（会话上下文聚合）}

\begin{lstlisting}[language=C++,caption={ShellState（节选）}]
struct ShellState {
    newdb::Session session;
    std::optional<newdb::Session::HeapAccess> session_heap_guard;
    TxnCoordinator txn;
    std::string log_file_path{"demo_log.bin"};
    std::string data_dir;
    bool encrypt_log{false};
    bool verbose{false};
    WhereQueryContext where_ctx;
};
\end{lstlisting}

说明：
\begin{itemize}
  \item \texttt{session\_heap\_guard} 是核心：在单条命令执行期内持有 \texttt{Session::mut\_}，保证缓存表指针有效；
  \item \texttt{where\_ctx} 保存查询缓存/LRU/策略状态/统计，避免每次 WHERE 都“冷启动”。
\end{itemize}

\subsubsection{DemoCliInvocation（一次 CLI 调用的完整语义）}

\begin{lstlisting}[language=C++,caption={DemoCliInvocation（节选）}]
struct DemoCliInvocation {
    DemoCliWorkspace ws;
    bool encrypt_log{false};
    bool verbose{false};
    DemoPrimaryAction primary{CliInteractive{}};
    std::string error;
};
\end{lstlisting}

它把“参数解析结果”标准化为 \texttt{workspace + flags + exactly one primary action}，后续 runner 可按 variant 分发。

\subsubsection{WhereQueryContext（查询缓存与策略状态）}

\begin{lstlisting}[language=C++,caption={WhereQueryContext（节选）}]
struct WhereQueryContext {
    std::unordered_map<std::string, std::vector<std::size_t>> query_cache;
    std::list<std::string> query_cache_lru;
    std::atomic<std::uint64_t> cache_hits{0};
    std::atomic<std::uint64_t> policy_rejects{0};
    WherePolicyState policy{};
    mutable std::mutex mu;
};
\end{lstlisting}

其作用是把“WHERE 的性能状态”与“策略限流状态”显式挂在会话层，而不是分散在静态全局变量中。

\subsection{核心流程}

\subsubsection{A. 入口主流程（demo\_main）}
\begin{enumerate}
  \item 解析 CLI：\texttt{demo\_parse\_invocation(argc, argv, inv)}
  \item 初始化运行上下文：\texttt{demo\_init\_session\_logging(app, ...)}
  \item 执行终端 phase（exec/batch/import/mdb）：
    \texttt{demo\_try\_run\_terminal\_phase(...)}
  \item 若无提前结束，进入交互 shell。
\end{enumerate}

\subsubsection{B. 命令处理主链（process\_command\_line）}
\begin{lstlisting}[language=C++,caption={命令分发主链（节选）}]
bool process_command_line(ShellState& st, const char* input_line) {
    // trim + logging + session commands
    if (handle_txn_commands(...)) return true;
    if (handle_workspace_admin_commands(...)) return true;
    if (handle_import_defattr_commands(...)) return true;
    if (handle_schema_catalog_commands(...)) return true;
    if (handle_ddl_create_use_commands(...)) return true;
    if (handle_ddl_alter_rename_commands(...)) return true;
    if (handle_schema_show_commands(...)) return true;
    if (handle_dml_insert_command(...)) return true;
    // fallback to cached HeapTable path
}
\end{lstlisting}

这条链的特点是“先高层语义命令、后表级缓存路径”，减少不必要的 heap 装载与解码。

\subsection{事务与回退能力（重构后现状）}

围绕事务路径，\fpath{cli/modules/txn/coordinator/*} 与命令分发层已形成“可执行 + 可恢复 + 可观测”的闭环：

\begin{itemize}
  \item \textbf{事务语法统一}：\texttt{SAVEPOINT / ROLLBACK TO / RELEASE SAVEPOINT} 在会话命令入口统一解释，降低命令分叉与 unknown 噪声；
  \item \textbf{任意回退语义}：支持按 savepoint 粒度回退与重做，命令层可组合为多步事务修复流；
  \item \textbf{崩溃恢复联动}：事务协调器与 WAL 的 analysis/redo/undo 恢复流程对齐，避免“命令层成功但恢复层不一致”；
  \item \textbf{诊断输出增强}：事务失败路径保留 soft/hard fail 诊断信息，便于自动化 gate 与 GUI 回看定位。
\end{itemize}

\subsection{算法与复杂度（中层视角）}

\begin{itemize}
  \item \textbf{CLI 分发}：参数扫描主要是 \(O(argc)\)；
  \item \textbf{命令路由}：按 handler 顺序匹配，最坏近似 \(O(H)\)（H=handler 数）；
  \item \textbf{WHERE}：
    \begin{itemize}
      \item id/索引命中接近 \(O(1)\) 或 \(O(\log n)\)（取决于底层结构）；
      \item fallback 全扫描路径为 \(O(n)\)；
      \item 缓存命中后可将重复查询候选构建降到接近 \(O(1)\)（不计输出）。
    \end{itemize}
\end{itemize}

\subsection{本章复评结论（现状 / 风险 / 建议）}

\begin{itemize}
  \item \textbf{分发结构已分层}：\texttt{registry/}+\texttt{router/}+\texttt{handlers/} 已落地，phase-2 前缀快路径缓解长链顺序匹配成本；后续演进重点转为“各 handler 内部复杂度”与“注册表生成源（快照 cc）体积治理”。
  \item \textbf{会话锁持有时长}：\texttt{session\_heap\_guard} 设计正确，但需警惕重 I/O 命令在持锁期间阻塞。
  \item \textbf{策略可观测性}：\texttt{WhereQueryContext} 计数与 \texttt{SHOW TUNING JSON}、\texttt{SHOW PLAN} 已部分贯通；可继续把缓存命中/拒绝原因结构化进 JSON，减少纯文本解析。
\end{itemize}

\subsection{近期修复要点：PAGE id 排序语义一致性}

在 \fpath{cli/modules/sidecar/page/page_index_sidecar.cc} 的近期修复中，\texttt{PAGE(..., id, asc)} 明确改为\textbf{按整数 id 比较}，避免字符串字典序导致的
\texttt{1,10,2} 问题（表现为新增到 10 后显示在 1 与 2 之间）。

\begin{itemize}
  \item \textbf{修复前}：\texttt{id} 在 sidecar 构建阶段以字符串键进入排序，可能出现 \texttt{10 < 2} 的错误顺序。
  \item \textbf{修复后}：在排序收敛阶段对 \texttt{order\_key=="id"} 走数值比较（\texttt{a.id < b.id} / \texttt{a.id > b.id}）。
  \item \textbf{兼容策略}：page sidecar header 引入版本字段（ver），旧索引自动失效重建，避免历史缓存继续污染分页顺序。
\end{itemize}

这类问题本质上是“\textbf{存储键编码语义} 与 \textbf{业务排序语义} 不一致”，后续新增索引 sidecar 时应优先遵循：
\begin{center}
排序比较规则必须与 \texttt{TableSchema::compare\_attr} / 字段类型语义保持一致。
\end{center}

